% =====================================================================
\section{Galois hull structure and main algebraic results}
\label{sec:hulls}
% =====================================================================

\begin{definition}\label{def:hull}
For $C\le\Rn$ and $s\ge0$, the \emph{$s$-Galois hull} of $C$ is
\[
  \Hull_s(C)=C\cap C^{\perp_s}.
\]
We write $h_s(C)=\dim_{\F_q}\Hull_s(C)$.
\end{definition}

\subsection{Structure of the hull}

\begin{theorem}[hull structure]\label{thm:hullstruct}
Let $C\le\Rn$ be a right $\Rn$-submodule. Then:
\begin{enumerate}[label=\textnormal{(\alph*)},leftmargin=*]
\item $\Hull_s(C)$ is a right $\Rn$-submodule;
\item $\Hull_s(C)$ is invariant under the cyclic shift $T$, i.e.
  $T\bigl(\Hull_s(C)\bigr)=\Hull_s(C)$; it is a right $\Rn$-submodule but
  \textbf{not} in general \emph{principal} --- that is
  Theorem~\ref{thm:nonprincipal} below;
\item $\Hull_s(C)$ is \textbf{not} in general a left ideal, hence not a
  two-sided ideal;
\item $\Hull_s(\sigma_kC)=\sigma_k\Hull_s(C)$ for every $k$, and
  $\Hull_s(\sigma_{2s}C)=\sigma_{2s}\Hull_s(C)$;
\item in the reflexive regime ($r\mid 2s$), $\Hull_s(C^{\perp_s})
  =\Hull_s(C)$ as subspaces; in general
  \[
    \Hull_s(C^{\perp_s})=C^{\perp_s}\cap\sigma_{2s}(C).
  \]
\end{enumerate}
\end{theorem}

\begin{proof}
(a) Intersection of two right $\Rn$-submodules
(Theorem~\ref{thm:duality}(a)); note that no principality hypothesis is used.

(b) $T$-invariance. Since $T$ and $\theta_s$ both act entrywise on the
coefficient tuple (the one by a cyclic index shift, the other by transpose and
Frobenius), they commute: $T\theta_s=\theta_sT$. The pairing
$\bb_s(x,y)=\Lambda\bigl(\sum_e\theta_s(x_e)y_e\bigr)$ is a sum over
positions, so $\bb_s(Tx,Ty)=\bb_s(x,y)$. Now $C$ is $T$-invariant and $T$ is a
bijection, so for $y\in R^n$,
\[
  \begin{aligned}
    y\in C^{\perp_s}
    &\iff \bb_s(x,y)=0\ \forall x\in C\\
    &\iff \bb_s(Tx,Ty)=0\ \forall x\in C\\
    &\iff \bb_s(x',Ty)=0\ \forall x'\in C\\
    &\iff Ty\in C^{\perp_s}.
  \end{aligned}
\]
Hence $C^{\perp_s}$ is $T$-invariant, and $T$ preserves the intersection
$C\cap C^{\perp_s}$. The \emph{principality} of the hull is not asserted here
and is false in general (Theorem~\ref{thm:nonprincipal}); what (a) and (b)
give is that $\Hull_s(C)$ is a right $\Rn$-submodule, which is all that
Section~\ref{sec:gray} and Section~\ref{sec:quantum} use.

(c) An explicit witness: at $(q,t,n)=(4,1,1)$, $s=0$, with
$g=\begin{pmatrix}0&1\\0&1\end{pmatrix}$, the hull $\Hull_0(g\Rn)$ is a right
ideal that is not closed under left multiplication by $R$.

(d) $\sigma_k$ is a ring automorphism fixing $X$, so it preserves right ideals
and commutes with $\theta_s$ up to the index shift
$\sigma_k\theta_s=\theta_s\sigma_k$; hence it carries $C^{\perp_s}$ to
$(\sigma_kC)^{\perp_s}$ and therefore the intersection to the intersection. The
second identity is the case $k$ replaced by $2s$, using
Theorem~\ref{thm:duality}(c).

(e) By Theorem~\ref{thm:duality}(c),
$(C^{\perp_s})^{\perp_s}=\sigma_{2s}(C)$, so
$\Hull_s(C^{\perp_s})=C^{\perp_s}\cap\sigma_{2s}(C)$. In the reflexive regime
$\sigma_{2s}=\id$, giving $\Hull_s(C^{\perp_s})=C^{\perp_s}\cap C=\Hull_s(C)$.
\end{proof}

\begin{proposition}[Morita hull formula and the principality criterion]
\label{prop:moritahull}
Under \eqref{eq:morita}, let $W=W(C)\le\Bq^2$ be the $\Bq$-submodule
corresponding to $C$ and $W'=W(C^{\perp_s})$. Then
\[
  \Hull_s(C)=I_{W\cap W'},
\]
the right ideal of $\Rn$ corresponding to $W\cap W'$, and
$\Hull_s(C)$ is principal if and only if $\MU_{\Bq}(W\cap W')\le2$.
\end{proposition}

\begin{proof}
Under the Morita correspondence, intersection of right ideals corresponds to
intersection of the associated submodules of $\Bq^2$, which gives the displayed
formula. The principality criterion is the Morita dictionary applied to the
minimal number of generators: a right ideal $I_U$ is principal iff $U$ is
generated by at most $2$ elements as a $\Bq$-module (the two standard basis
vectors of $\Bq^2$ being the generators of $\Rn$ as a right module).
\end{proof}

\subsection{Non-principal hulls: the conjecture O-3.1 is false}

Over a finite chain ring every submodule of $S^2$ is free of rank at most $2$,
hence $2$-generated. Combined with Lemma~\ref{lem:components} this suggests the
following natural conjecture, which was the open problem O-3.1 of this project.

\begin{quote}\textbf{Conjecture O-3.1 (false, as stated).} For every cyclic
code $C$ over $R=M_2(A_t)$ and every $s$, the Galois hull $\Hull_s(C)$ is a
principal right ideal of $\Rn$.
\end{quote}

Two readings must be separated, because only the first is refuted here:
\begin{description}
\item[(O-3.1a)] for \emph{every} cyclic code $C$ (i.e. every right ideal of
  $\Rn$, principal or not) the hull is principal. --- \textbf{Refuted},
  Theorem~\ref{thm:nonprincipal}.
\item[(O-3.1b)] for every \emph{principal} $C=g\Rn$ the hull is principal.
  --- \textbf{Open}; see Remark~\ref{rem:o31b}.
\end{description}

\begin{theorem}\label{thm:nonprincipal}
Conjecture O-3.1 is false. Moreover:
\begin{enumerate}[label=\textnormal{(\alph*)},leftmargin=*]
\item $\MU_{\Bq}(\Hull_s(C))\ge3$ requires $t\ge2$ \textbf{and} $p\mid n$;
\item the bound $\MU_{\Bq}(W)\le2$ is false over a non-chain local
  component, and it fails for a hull.
\end{enumerate}
\end{theorem}

\begin{proof}
(a) By Lemma~\ref{lem:components}, $\MU_{\Bq}(W)=\max_j\MU_{\Bq_j}(W_j)$, and
$\Bq_j$ is a chain ring iff $t=1$ or $e_j=1$. If $t=1$ then every component is a
chain ring and every $W_j$ is $2$-generated, so $\MU_{\Bq}(W)\le2$. If
$p\nmid n$ then $X^n-1$ is separable over $\F_q$, so every $e_j=1$ and again
every component is a chain ring. Hence $\MU_{\Bq}(W)\ge3$ forces $t\ge2$ and
$p\mid n$ (equivalently: some $e_j\ge2$, which happens exactly when $p\mid n$,
since $X^n-1=(X^{n'}-1)^{p^v}$ with $n=p^vn'$ and $p\nmid n'$).

(b) Let $\Bq_j\cong\F_q[u,w]/\langle u^2,w^2\rangle$ be a wild component. Its
maximal ideal is $\mathfrak m=(u,w)$, $\mathfrak m^2=(uw)$,
$\mathfrak m^3=0$, and $\ann_{\Bq_j^2}(\mathfrak m)$ has $\F_q$-dimension $2$.
Consider the family of submodules
\begin{equation}\label{eq:witnessfamily}
  W=\mathfrak m^ie_1\oplus U,\qquad U\le\ann_{\Bq^2}(\mathfrak m),
\end{equation}
where $e_1$ is the first standard basis vector of $\Bq^2$. Then
$\mathfrak mW=\mathfrak m^{i+1}e_1$, so by Nakayama
\[
  \MU_{\Bq}(W)=\dim_{\F_q}\mathfrak m^i-\dim_{\F_q}\mathfrak m^{i+1}
  +\dim_{\F_q}U .
\]
Taking $i=0$ (so $\dim\mathfrak m^0-\dim\mathfrak m^1=\dim\Bq-\dim\mathfrak m
=4-3=1$ over $\F_q$ for $t=2$, $e_j=2$) and $\dim U=2$ gives
$\MU_{\Bq}(W)=3$. Enumerating the deterministic family \eqref{eq:witnessfamily}
exhaustively in the wild families produces submodules whose Galois hull has
$\MU_{\Bq}=3$; see Section~\ref{subsec:witnesses} for the explicit witnesses.
\end{proof}

\begin{remark}
The family \eqref{eq:witnessfamily} is finite and completely deterministic, so
the witnesses are reproducible without search. The use of $e_1$ and of
$\ann(\mathfrak m)$ is essential: the attempted construction
\[
  W=\mathfrak m^ie_1\oplus U,\qquad U\le\Bq^2\ \text{arbitrary},
\]
does \emph{not} give a $\Bq$-submodule, and the direct
``complement-lift'' construction fails.
\end{remark}

\subsection{The witnesses}\label{subsec:witnesses}

The following families have been examined; in each, the deterministic family
\eqref{eq:witnessfamily} was enumerated completely and the Galois hull
$\Hull_s(W)$ computed by the Morita route, with $\MU_{\Bq}$ of the hull
recomputed directly (not inferred).

\begin{table}[ht]
\centering
\caption{Wild families: ideal structure and the best witness found. Here
$\dim\mathfrak m^i$ means $\dim_{\F_q}\mathfrak m^i$.}
\label{tab:witnesses}
\small
\begin{tabular}{@{}ccccccc@{}}
\toprule
$(q,t,n,s)$ & $\Nlen$ & $\dim\mathfrak m,\dim\mathfrak m^2,\dots$
 & $\dim\ann(\mathfrak m)$ & $\#\{W:\MU_{\Bq}(W)\ge3\}$
 & $\MU_{\Bq}(\Hull_s)$ & $\Hull_s=W$? \\
\midrule
$(4,2,2,0)$ & 16 & $3,1$ & 2 & 12 & 2 & --- \\
$(4,2,2,1)$ & 16 & $3,1$ & 2 & 12 & 2 & --- \\
$(8,2,2,0)$ & 16 & $3,1$ & 2 & 56 & 2 & --- \\
$(8,2,2,1)$ & 16 & $3,1$ & 2 & 56 & 2 & --- \\
$(9,2,3,0)$ & 24 & $5,3,1$ & 2 & 144 & \textbf{3} & yes \\
$(9,2,3,1)$ & 24 & $5,3,1$ & 2 & 144 & \textbf{3} & yes \\
$(4,2,4,0)$ & 32 & $7,5,3,1$ & 2 & 36 & \textbf{3} & no \\
$(4,2,4,1)$ & 32 & $7,5,3,1$ & 2 & 36 & \textbf{3} & no \\
\bottomrule
\end{tabular}
\end{table}

In addition, an exhaustive enumeration of \emph{every} code of \emph{every}
enumerable family (families of sizes $7,33,33,1221,343,144,93,35937$) found
four hulls with $\MU_{\Bq}=3$ at $(4,2,2)$ --- one at $s=0$ and three at $s=1$
--- which is the original refutation of O-3.1.

\begin{remark}[the sharper question (O-3.1b) remains open]\label{rem:o31b}
Every witness found --- the four at $(4,2,2)$ and the four at $(9,2,3)$ and
$(4,2,4)$ --- has $\MU_{\Bq}(W)=3$ for the \emph{code} as well as for the
hull. Equivalently, every known witness is a code that is itself
non-principal. The four $(4,2,2)$ witnesses and the two $(9,2,3)$ witnesses
additionally satisfy $\Hull_s(W)=W$, so the hull is non-principal because the
code is.

It is therefore \textbf{open} whether a \emph{principal} cyclic code can have a
non-principal Galois hull in the wild regime. We do not claim it either way,
and we note that the two questions are frequently conflated. What
Theorem~\ref{thm:nonprincipal} settles is that hull principality cannot be
assumed, so no downstream argument may use it.
\end{remark}

\begin{remark}[a disambiguation]\label{rem:mub}
The number $\MU_B$ is attached to a \emph{module}. The value $\MU_B=4$ reported
in earlier work of this project refers to $\MU_B(W)$ for the \emph{code}
$W$ at $(4,2,2,1)$; the value constrained by the hull census is
$\MU_B(\Hull_s(W))$, whose observed range is $\{0,1,2,3\}$ and whose maximum
over the extended enumeration is $3$. The two numbers are not in conflict.
\end{remark}

\subsection{What does not exist}

\begin{theorem}[no closed hull formula]\label{thm:noclosed}
There is no universal closed formula of ``gcd type'' for the Galois hull of a
cyclic code over $R=M_2(A_t)$: that is, there is no expression in the
generators of $C$, built from the ring operations and a finite fixed set of
polynomial operations, that yields a generator of $\Hull_s(C)$ for all $C$ and
all $s$.
\end{theorem}

\begin{proof}
If such a formula existed, every hull would be principal, since the output of
the formula is (by construction) a single element generating the hull. This
contradicts Theorem~\ref{thm:nonprincipal}.
\end{proof}

\begin{remark}
One should not read Theorem~\ref{thm:noclosed} as a negative result about the
\emph{usefulness} of hulls. The Morita formula of
Proposition~\ref{prop:moritahull} is a perfectly effective way to compute
$\Hull_s(C)$: it is a submodule intersection in $\Bq^2$, and via the CRT
decomposition \eqref{eq:crt} it is computed componentwise. What does not exist
is a generator-polynomial formula of the familiar commutative kind.
\end{remark}

\subsection{LCD, self-orthogonal and self-dual codes}

\begin{proposition}\label{prop:lcdso}
Let $C\le\Rn$ with Morita module $W=W(C)$ and $W'=W(C^{\perp_s})$. Then
\begin{enumerate}[label=\textnormal{(\alph*)},leftmargin=*]
\item $C$ is LCD ($\Hull_s(C)=\{0\}$) iff $W\cap W'=\{0\}$ iff
  $C\oplus C^{\perp_s}=R^n$;
\item $C\subseteq C^{\perp_s}$ iff $W\subseteq W'$; for $C=g\Rn$ principal
  this is equivalent to $\theta_s(\rev g)\,g=0$;
\item $C=C^{\perp_s}$ iff $C\subseteq C^{\perp_s}$ and $2\dim_{\F_q}C=\Nlen$;
  for $C=g\Rn$ this is equivalent to $\im g=\ker\theta_s(\rev g)$.
\end{enumerate}
\end{proposition}

\begin{proof}
Immediate from Proposition~\ref{prop:moritahull} and the dimension count of
Theorem~\ref{thm:duality}(b); the principal-case criteria follow by translating
$W=\im g$ and $W'=\ker\theta_s(\rev g)$ through
Theorem~\ref{thm:dualformula}.
\end{proof}

\begin{remark}[two conjectures, stated as such]
Let $D=C^{\perp_s}$. It is natural to ask whether
$\Hull_s(C)\subseteq\sigma_{2s}(C)$ always holds, equivalently whether
$\Hull_s(D)=\Hull_s(D^{\perp_s})$; and whether
$\dim\Hull_s(C)=\dim\Hull_s(C^{\perp_s})$ always. Both hold in every instance
computed, and the first would follow from $\Hull_s(C^{\perp_s})\subseteq
\sigma_{2s}(C^{\perp_s})$. We state them as conjectures: neither is used
anywhere in this paper, and neither is claimed as a theorem.
\end{remark}

\subsection{Worked examples III: hulls}

\begin{example}[an LCD code, by the Morita criterion]\label{ex:lcd}
With the notation of Example~\ref{ex:dual} --- $(q,t,n)=(4,1,1)$, $s=0$ --- the
code $C=I_W$ with $W=\F_4(1,0)^{\mathsf T}$ has $W\cap W^{\perp_0}=\{0\}$, so
by Proposition~\ref{prop:moritahull}
\[
  \Hull_0(C)=I_{\{0\}}=\{0\},\qquad h_0(C)=0,
\]
and $C\oplus C^{\perp_0}=\Rn$; $C$ is LCD. The other extreme is $W=\F_4^2$, for
which $C=\Rn$ and $C^{\perp_0}=\{0\}$, again LCD.

The complete picture in this family is small enough to state: $\F_4^2$ has
$1+5+1=7$ subspaces, and a $1$-dimensional $W=\F_4v$ is self-orthogonal iff
$v\cdot v=0$. Since $v\cdot v=v_0^2+v_1^2=(v_0+v_1)^2$ in characteristic $2$,
the self-orthogonal lines are exactly those with $v_0+v_1=0$, of which there is
one, namely $v=(1,1)$. So of the $7$ codes, $6$ are LCD and one is self-dual.
\end{example}

\begin{example}[a self-orthogonal and a self-dual code]\label{ex:sd}
Take $W=\F_4\cdot(1,\omega)^{\mathsf T}$ with $\omega\in\F_4$ a root of
$\omega^2+\omega+1=0$, again at $(q,t,n)=(4,1,1)$ and $s=0$. Then
$v\cdot v=1+\omega^2=1+(\omega+1)=\omega\ne0$, so $W\nsubseteq W^{\perp_0}$
and $C=I_W$ is LCD.

Taking instead $W=\F_4\cdot(1,1)^{\mathsf T}$ gives $v\cdot v=1+1=0$, so
$W\subseteq W^{\perp_0}$. Both spaces have dimension $1$, hence
$W=W^{\perp_0}$, and therefore
\[
  C=C^{\perp_0},\qquad \Hull_0(C)=C,\qquad h_0(C)=\dim C=2=\Nlen/2 :
\]
the code is self-dual. By Proposition~\ref{prop:lcdso}(c), for a principal
$C=g\Rn$ the criterion is $\im g=\ker\theta_s(\rev g)$, which in this family
(where $\rev=\id$ by Corollary~\ref{cor:rev}) is a $4\times4$ rank
computation over $\F_4$.
\end{example}

\begin{example}[the $(9,2,3)$ witness, in full]\label{ex:witnessA}
Let $q=9$, $t=2$, $n=3$. Then $X^3-1=(X-1)^3$ over $\F_9$, so
\[
  \Bq\cong\F_9[u,w]/\langle u^2,w^3\rangle ,
\]
local with maximal ideal $\mathfrak m=(u,w)$ satisfying
\[
  \dim_{\F_9}\mathfrak m=5,\quad \dim_{\F_9}\mathfrak m^2=3,\quad
  \dim_{\F_9}\mathfrak m^3=1,\qquad \dim_{\F_9}\ann_{\Bq^2}(\mathfrak m)=2 .
\]
This is a wild component ($t\ge2$ and $e_1=3\ge2$), so by
Lemma~\ref{lem:components} the bound $\MU_{\Bq}\le2$ is not available here.
Enumerate the deterministic family \eqref{eq:witnessfamily},
\[
  W=\mathfrak m^ie_1\oplus U,\qquad U\le\ann_{\Bq^2}(\mathfrak m),
\]
completely: $144$ of its members have $\MU_{\Bq}(W)\ge3$. For each, form the
Galois hull $\Hull_s(W)$ by the Morita route and recompute
$\MU_{\Bq}(\Hull_s(W))$ directly. The outcome for $s=0$ and for $s=1$ is
\[
  \MU_{\Bq}(W)=3,\qquad \MU_{\Bq}\bigl(\Hull_s(W)\bigr)=3,\qquad
  \Hull_s(W)=W .
\]
The last equality says the code equals its own Galois hull. The corresponding
image code $D=\varphi_3(C)\subseteq\F_9^{24}$ has $k=\dim D=10$ and Euclidean
hull $h_E=8$, so route R3 gives
\[
  [[24,\,24-10-8,\,2\,;\,10-8]]_9=[[24,6,2;2]]_9 ,
\]
while route R1 gives $[[24,14,2]]_9$ with an abelian stabilizer group. The
parameters are weak; the point is that the hull is \emph{not principal}, and
yet every step of the construction applies to it unchanged. Together with
Example~\ref{ex:witnessB} this is the refutation of O-3.1 outside the
originally known family $(4,2,2)$.
\end{example}

\begin{example}[the $(4,2,4)$ witness]\label{ex:witnessB}
Let $q=4$, $t=2$, $n=4$. Then $X^4-1=(X-1)^4$ over $\F_4$ and
\[
  \begin{gathered}
    \Bq\cong\F_4[u,w]/\langle u^2,w^4\rangle,\\
    \dim_{\F_4}\mathfrak m=7,\ \dim\mathfrak m^2=5,\
      \dim\mathfrak m^3=3,\ \dim\mathfrak m^4=1,\\
    \dim_{\F_4}\ann_{\Bq^2}(\mathfrak m)=2 .
  \end{gathered}
\]
The same construction \eqref{eq:witnessfamily} has $36$ members with
$\MU_{\Bq}(W)\ge3$, and for $s=0$ and $s=1$ the best witness satisfies
\[
  \MU_{\Bq}(W)=3,\qquad \MU_{\Bq}\bigl(\Hull_s(W)\bigr)=3,\qquad
  \Hull_s(W)\ne W .
\]
Here $\Nlen=4tn=32$, $k=14$, $h_E=10$, so route R3 gives
\[
  [[32,\,32-14-10,\,2\,;\,14-10]]_4=[[32,8,2;4]]_4 .
\]
Unlike Example~\ref{ex:witnessA}, the hull is a proper superobject of the
code one started from. It is nonetheless carried unchanged through the Gray
map by Theorem~\ref{thm:grayhull}, which is precisely the statement that needs
no principality.
\end{example}
